problem:  
Let \((G, d_{\Delta})\) be a connected, simply connected, **equiregular Carnot group** endowed with its sub‑Riemannian (Carnot–Carathéodory) metric \(d_{\Delta}\) and Haar measure \(\mu\). For every integer \(k \ge 1\), define the **homogeneous filling function**  
\[
\mathrm{Fill}_{G}^{(k)}(r)=\sup\left\{ \inf\{ \mu_{k+1}(T)\mid \partial T=S \},\ \mu_k(S)\le r \right\},
\]  
where \(S\) ranges over \(k\)-dimensional integral currents of compact support in \(G\) with \(\partial S=0\), and \(\mu_k\) denotes the sub‑Riemannian \(k\)-dimensional Hausdorff measure.  

It is known in special cases (e.g., Heisenberg or abelian groups) that there exists an exponent \(\beta(k,G)\) such that \(\mathrm{Fill}_G^{(k)}(r) \simeq r^{\beta(k,G)}\).  

**Problem:**  
For general equiregular Carnot groups \(G\) of homogeneous (Hausdorff) dimension \(Q\), determine whether there exists a **universal structural formula** for the filling exponent  
\[
\beta(k,G)
\]  
expressible purely in terms of the **growth vector** of \(G\)—that is, in terms of the sequence \((m_1,\dots,m_s)\) where \(m_i = \dim(V_1\oplus\dots\oplus V_i)\) and \(s\) is the step of \(G\).  

Equivalently:  
Is there an explicit function \(F: \mathbb{N}^s \times \mathbb{N}\to \mathbb{R}^+\) such that for every equiregular Carnot group \(G\) with growth vector \((m_1,\dots,m_s)\) and for all \(k\),
\[
\beta(k,G)=F(m_1,\dots,m_s;k)
\]
and conversely—if such a function exists—determine whether it detects the differentiable (or metric) equivalence class of \(G\).  

Why is it a "good" problem:  
This problem refines and extends Gromov’s notion of filling invariants to the precise sub‑Riemannian realm. It asks whether the asymptotic geometric complexity of a Carnot group, captured by its filling exponents, can be expressed purely in algebraic (Lie–theoretic) data—the growth vector.  

It is a good problem because:  
- **Profound insight:** It seeks a universal structural law connecting large‑scale geometric measure invariants (filling functions) with microscopic algebraic data (the Lie algebra stratification). Resolving it would give a geometric–analytic fingerprint of sub‑Riemannian spaces.  
- **Bridge between fields:** It intertwines sub‑Riemannian geometry, geometric group theory (as Carnot groups are nilpotent Lie groups), and geometric measure theory (via integral currents). Any progress would likely require developing new analytic or algebraic tools bridging these domains.  
- **Extreme simplicity:** The question can be stated in one line—*Is the filling exponent determined purely by the growth vector?*—yet the implications are profound, potentially leading to a new classification scheme for Carnot geometries and shedding light on how measure‑theoretic filling behavior encodes algebraic structure.